% =============================================================================
% Chapter 9: Method -- Part II
% =============================================================================
\mychapter{9. Method}

\mysection{9.1 PODFAM Sample Set}
Five XCT-scanned metal AM samples were used for the core validation: sample\_01, sample\_02,
sample\_03, sample\_04, and sample\_05, comprising 900 slices each (sample\_04: 749 slices) at
approximately $1010\times980$--988 pixel resolution, 16-bit unsigned TIFF. As with the Part I
dataset, no expert-annotated ground truth is available for any PODFAM sample. Section~10.5
additionally reports on sample\_07 (1,473 slices, $1143\times1215$ pixels), acquired and processed
after this core validation.

All PODFAM samples are Ti-6Al-4V, manufactured by laser powder bed fusion (PBF-LB) on an Aconity Two
machine. Post-build, the samples were characterised by X-ray computed tomography on a Waygate
v\textbar tome\textbar x M~240 XCT system. Material and process details were not part of the original
mid-term submission and are added here as they were not required for, and do not affect, the
preprocessing, pseudo-labelling, or detector-training methodology described in Sections~9.2--9.7,
which operate purely on the reconstructed greyscale volumes; they are recorded here for completeness
and for anyone reproducing this validation on a different Ti-6Al-4V PBF-LB dataset to judge
comparability.

\mysection{9.2 Preprocessing}
Each slice is processed independently by (i) a $3\times3$ median filter to suppress speckle noise,
followed by (ii) non-local means (NLM) denoising with adaptive filter strength
$h = 0.6 \times \hat{\sigma}$, where $\hat{\sigma}$ is the per-slice estimated noise standard
deviation, and (iii) normalisation to an 8-bit intensity range. This is a narrower pipeline than the
four-stage Part I pipeline: beam-hardening correction and polar-coordinate ring-artefact suppression
were not applied to the PODFAM data in this validation phase, since PODFAM's images did not exhibit
those particular artefacts as strongly and the analysis prioritised comparability across the
classical-thresholding and U-Net stages over full artefact removal.

\mysection{9.3 Sample-Boundary Detection and Pseudo-Label Generation}
A circular sample-boundary mask is fitted once per volume by averaging the detected circle centre and
radius across five evenly-spaced slices, then eroding the resulting disc by a small margin (5~px) to
exclude partial-volume edge pixels, following the same rationale as the Part I standard-method
erosion step~\citep{NIST_Standard}. All three thresholding methods (Otsu, Yen, Bernsen) are then
computed only within this interior region, preventing background air from being misclassified as
defect. Bernsen's local contrast threshold (DCT) is computed automatically per image from the
estimated local noise standard deviation of the solid phase, following Kim et
al.~\citep{Kim2017}; Bernsen is used as the primary pseudo-label source for U-Net training, motivated
directly by the comparative finding reported in Section~8.1.

\mysection{9.4 Defect-Content-Ranked Sample Selection}
Rather than assigning samples to train/validation/test roles alphabetically or at random, every
sample's mean Bernsen defect fraction is computed over its full volume and samples are ranked
most-voids-first. The three samples retained for this validation phase were, in ranked order,
sample\_02 (7.586\% defect fraction), sample\_04 (1.994\%), and sample\_01 (0.043\%); sample\_03 was
excluded as near-empty (0.0004\%, effectively uninformative for training), and sample\_05 was
excluded for the reason given below.

The most defect-rich sample was assigned to training, the next to validation, and the least to test,
deterministically, rather than via a random index-based split. This choice is methodologically
significant: an earlier iteration of this pipeline used a fixed-seed random split over the same three
ranked samples, which by chance assigned the single most defect-rich sample to the test role rather
than training. The resulting model, trained on a far less informative sample, could not learn the
defect pattern that was concentrated almost entirely in the sample it never saw, producing a held-out
Dice score below 0.15 for reasons unrelated to model capacity or training procedure. This is reported
here as a methodological finding in its own right: with only three usable samples, split assignment
must be deliberate rather than left to a fixed random seed.

\mysection{9.5 A Documented Pseudo-Label Failure Mode}
During full-volume evaluation (Section~10.1), sample\_05 was found to have an implausible Bernsen
defect fraction of 52.8\%, physically inconsistent with a metal AM part, and roughly three orders of
magnitude above every other sample's Bernsen result. Investigation traced this to Bernsen's automatic
contrast-threshold (DCT) estimation: for this dataset, the estimated DCT (329--520) is far above the
value the underlying method's own reference reports as typical (approximately 15)~\citep{Kim2017}.
An inflated DCT routes almost every pixel into Bernsen's low-contrast fallback branch, which
classifies purely by a fixed intensity threshold at 128. Samples whose preprocessed images average
brighter than 128 are still classified approximately correctly by this fallback; sample\_05 averages
79.9, substantially darker, so the fallback misclassifies most of the genuinely solid material as
void. Otsu and Yen independently agree with Bernsen's implausible result on this specific sample
(51.7\% and 47.9\% respectively), which is consistent with a shared upstream cause (the sample's
unusually low mean intensity) rather than three independent confirmations of real porosity.
sample\_05 was excluded from training and from the ranked comparison via an automatic guard that
flags any sample exceeding a 10\% defect-fraction threshold, a level far above plausible AM porosity,
for manual review rather than allowing it to silently dominate a ranking or a training run.

\mysection{9.6 Detection Model and Training Configuration}
A 2D U-Net (approximately 31.0 million trainable parameters; four encoder/decoder stages, feature
channels $64\to128\to256\to512$, bottleneck dropout $p = 0.2$) is trained on $256\times256$ pixel
patches (stride 128, 50\% overlap) drawn from 80 evenly-spaced slices per training volume, using
stratified 1:3 foreground-to-background patch sampling to address the severe class imbalance
inherent to porosity data. Patches are shuffled at the slice level rather than individually during
training, all patches from a visited slice are drawn together before moving to the next, to keep
memory-mapped volume access efficient on constrained hardware without sacrificing epoch-level
randomisation. The combined Dice-Focal loss ($\lambda = 0.5$, $\alpha = 0.25$, $\gamma = 2.0$) is
used, with Adam optimisation (learning rate $4\times10^{-5}$), early stopping (patience 10 epochs),
and best-checkpoint selection by validation Dice.

\mysection{9.7 Evaluation Metrics}
Detection performance is reported using the Dice coefficient, Intersection-over-Union (IoU),
precision, and recall, computed at the pixel level against the Bernsen pseudo-label. Full-volume
defect statistics (defect fraction, pore count, mean pore area, mean equivalent diameter) are
computed by pooling per-slice connected-component analysis across every slice in a volume, rather
than from a single representative slice, for both the classical methods and the trained detector's
predictions, ensuring the classical and deep-learning results in Section~10 are directly comparable
rather than mixing single-slice and full-volume statistics.
